\section*{Capítulo \ref{ch:b3:seismology-imaging} --- Sismologia e imageamento da Terra}

\begin{solution}{exo:b3:seismology-imaging:1}
(a) $\sin i_0 = pv_0/R = 500\times5.80/6371 = 0.4552$, logo $i_0 = \ang{27.1}$.
(b) Em \qty{771}{km}, $r = \qty{5600}{km}$ e $r/v = 5600/11.066 =
\qty{506.1}{s.rad^{-1}}$; em \qty{871}{km}, $r = \qty{5500}{km}$ e $r/v =
5500/11.245 = \qty{489.1}{s.rad^{-1}}$. O raio atinge o ponto de retorno onde $r/v = p = 500$:
$z_t = 771 + 100\times(506.1 - 500)/(506.1 - 489.1) = 771 + 36 =
\qty{807}{km}$.
\end{solution}

\begin{solution}{exo:b3:seismology-imaging:2}
(a) $(455.5 - 369.6)/10 = \qty{8.59}{s}$ por grau, e
$8.59\times180/\pi = \qty{492.2}{s.rad^{-1}}$. (b) Um grau de arco equivale a
$6371\pi/180 = \qty{111.19}{km}$, logo a frente de onda varre a superfície a
$111.19/8.59 = \qty{12.94}{km.s^{-1}}$, muito mais rápido do que qualquer rocha na
superfície: o raio chega quase verticalmente. (c) No ponto de retorno $v = r_t/p =
5450/492.2 = \qty{11.07}{km.s^{-1}}$, uma velocidade da onda P no manto inferior.
\end{solution}

\begin{solution}{exo:b3:seismology-imaging:3}
(a) $\mu_1 = \rho_1\beta_1^2 = 2700\times3500^2 = \qty{3.31e10}{Pa}$
(\qty{33.1}{GPa}); $\mu_2 = 3300\times4500^2 = \qty{6.68e10}{Pa}$
(\qty{66.8}{GPa}). (b) $\lambda = cT = 3.79\times20 = \qty{75.8}{km}$, 2.17
vezes a espessura da camada: a onda alcança bem o interior do semiespaço, o que
explica por que sua velocidade já está acima de $\beta_1$. (c) $4000/3.384 = \qty{1182}{s}$ e
$4000/4.333 = \qty{923}{s}$: as ondas de 20 s chegam \qty{259}{s}, ou
\qty{4.3}{min}, depois das ondas de 100 s.
\end{solution}

\begin{solution}{exo:b3:seismology-imaging:4}
(a) ${}_0S_2$: $1/(53.89\times60) = \qty{3.093e-4}{Hz} = \qty{0.3093}{mHz}$;
${}_0T_2$: $1/(44.01\times60) = \qty{0.3787}{mHz}$. (b) $1440/53.89 = 26.7$
e $1440/44.01 = 32.7$ oscilações por dia. (c) A diferença é
$\qty{0.0694}{mHz}$, cujo inverso é \qty{14400}{s}: no mínimo \qty{4.0}{h} de
registro, na prática muito mais, já que os picos são alargados pela atenuação e
desdobrados.
\end{solution}

\begin{solution}{exo:b3:seismology-imaging:5}
(a) A corda entre dois pontos da esfera separados pelo ângulo $\Delta$
no centro tem comprimento $2R\sin(\Delta/2)$, percorrida à velocidade $v$.
(b) Para \ang{30}: $2\times6371\times\sin\ang{15}/10.0 = \qty{329.8}{s}$, contra
\qty{369.6}{s}; para \ang{90}: $2\times6371\times\sin\ang{45}/10.0 =
\qty{901.0}{s}$, contra \qty{779.7}{s}. (c) O raio para \ang{30} atinge seu ponto de retorno no
manto superior, onde a velocidade real é de cerca de 8 a \qty{10}{km.s^{-1}}, abaixo da
uniforme de \qty{10}{km.s^{-1}}; o raio para \ang{90} percorre a maior parte do trajeto no
manto inferior, a 11 a \qty{13.7}{km.s^{-1}}, e se curva para baixo ao entrar nele, logo a
Terra real é mais rápida.
\end{solution}

\begin{solution}{exo:b3:seismology-imaging:6}
(a) Elevando ao quadrado, $(2 - \xi)^4 = 16(1 - \xi/3)(1 - \xi)$. Desenvolvendo:
$16 - 32\xi + 24\xi^2 - 8\xi^3 + \xi^4 = 16 - \tfrac{64}{3}\xi + \tfrac{16}{3}\xi^2$,
logo $\xi^4 - 8\xi^3 + \tfrac{56}{3}\xi^2 - \tfrac{32}{3}\xi = 0$, isto é
$\xi\,(3\xi^3 - 24\xi^2 + 56\xi - 32) = 0$, e
$3\xi^3 - 24\xi^2 + 56\xi - 32 = (\xi - 4)(3\xi^2 - 12\xi + 8)$. A raiz
$\xi = 0$ não corresponde a uma onda. (b) As raízes da quadrática são $2 \pm 2/\sqrt3$,
isto é 3.155 e 0.8453; somente $\xi = 0.8453$ está em $(0, 1)$ (e $\xi = 4$
não está), e ele satisfaz a equação original, sem elevação ao quadrado, ambos os lados valendo $4/3$.
Logo $c/\beta = \sqrt{0.8453} = 0.9194$. (c) $c = 0.9194\times3.46 =
\qty{3.181}{km.s^{-1}}$; ao longo de \qty{1000}{km}, $314.4 - 289.0 = \qty{25.3}{s}$
depois da onda S.
\end{solution}

\begin{solution}{exo:b3:seismology-imaging:7}
$\mathsf G$ tem as linhas $(1, x_i)$ e $\mathsf C_d^{-1} = \operatorname{diag}
(100, 100, 25, 25)$ (\unit{s^{-2}}). As equações normais são
\[
  \begin{pmatrix} 250 & 23750 \\ 23750 & \num{2812500} \end{pmatrix}
  \begin{pmatrix} a \\ b \end{pmatrix}
  = \begin{pmatrix} 3665 \\ 418125 \end{pmatrix},
\]
com determinante \num{1.390625e8}. Logo $b = (250\times418125 -
23750\times3665)/\num{1.390625e8} = \qty{0.12575}{s.km^{-1}}$ e
$a = (2812500\times3665 - 23750\times418125)/\num{1.390625e8} =
\qty{2.713}{s}$. A covariância é $\mathsf N^{-1}$: $\sigma_b =
\sqrt{250/\num{1.390625e8}} = \qty{0.00134}{s.km^{-1}}$, $\sigma_a =
\sqrt{2812500/\num{1.390625e8}} = \qty{0.142}{s}$. A velocidade é
$1/b = \qty{7.95}{km.s^{-1}}$, com o erro $\sigma_b/b^2 =
\qty{0.08}{km.s^{-1}}$: uma velocidade do topo do manto. Os resíduos são $-0.001$,
$0.011$, $-0.076$ e $\qty{0.036}{s}$, e $\chi^2 = 0.19$ para $4 - 2 = 2$
graus de liberdade: menor do que o esperado, logo os erros informados são, se algo, generosos;
o ajuste é aceitável.
\end{solution}

\begin{solution}{exo:b3:seismology-imaging:8}
(a) $\mathsf G = (30\;\;40)$ é uma única linha: $\mathsf G^{\mathsf T}\mathsf G$ é
$2\times2$ de posto um, logo singular; qualquer modelo com $30m_1 + 40m_2 = 0$ pode
ser somado sem alterar o dado. (b) $\mathsf G\mathsf G^{\mathsf T} =
900 + 1600 = 2500$, logo $\mathbf m = (30, 40)\times0.50/2500 =
(0.006, 0.008)\,\unit{s.km^{-1}}$. (c) $\mathsf R = \frac{1}{2500}
\begin{pmatrix} 900 & 1200 \\ 1200 & 1600 \end{pmatrix} =
\begin{pmatrix} 0.36 & 0.48 \\ 0.48 & 0.64 \end{pmatrix}$, de traço 1: um
dado resolve uma combinação dos parâmetros. (d) $\mathsf R\,(0.01, 0) =
(0.0036, 0.0048)\,\unit{s.km^{-1}}$: a anomalia se distribui ao longo do raio, mais
na célula onde o raio é mais longo.
\end{solution}

\begin{solution}{exo:b3:seismology-imaging:9}
(a) Linearizado: $\delta t = -(\delta v/v)L/v = -0.03\times200/9.0 =
\qty{-0.667}{s}$; exatamente, $L/(1.03v) - L/v = 200/9.27 - 200/9.0 =
\qty{-0.647}{s}$. (b) O trajeto na laje é $100/\cos\ang{40} =
\qty{130.5}{km}$: $-0.03\times130.5/9.0 = \qty{-0.435}{s}$ linearizado e
$\qty{-0.422}{s}$ exatamente. (c) O valor exato é o linearizado dividido por
$1.03$: um erro de \qty{3}{\percent}, da ordem do próprio $\delta v/v$, como
se espera de uma teoria de primeira ordem.
\end{solution}

\begin{solution}{exo:b3:seismology-imaging:10}
(a) O sobretom existe enquanto sua parte no semiespaço decai, $c < \beta_2$; no
corte $c = \beta_2$, logo $q_2 = 0$ e $\tan(kHq_1) = 0$, com $kHq_1$ em
$[n\pi, n\pi + \pi/2)$: $kHq_1 = n\pi$. (b) Com $k = \omega/\beta_2$ e
$q_1 = \sqrt{\beta_2^2/\beta_1^2 - 1} = \sqrt{\beta_2^2 - \beta_1^2}/\beta_1$,
$\omega_nH\sqrt{\beta_2^2 - \beta_1^2}/(\beta_1\beta_2) = n\pi$, que é a
fórmula. (c) $\sqrt{4.5^2 - 3.5^2} = \sqrt8 = 2.828$; $\omega_1 =
\pi\times3.5\times4.5/(35\times2.828) = \qty{0.4998}{rad.s^{-1}}$, um período de
$2\pi/0.4998 = \qty{12.57}{s}$: em períodos mais longos apenas o modo fundamental
existe. (d) Para $n = 0$ a frequência de corte é zero: o modo fundamental
existe em todos os períodos. Em períodos curtos $kH \to \infty$ enquanto $kHq_1$ permanece
abaixo de $\pi/2$, logo $q_1 \to 0$ e $c \to \beta_1$.
\end{solution}

\begin{solution}{exo:b3:seismology-imaging:11}
(a) $\mathsf G^{\mathsf T}(\mathsf G\mathsf G^{\mathsf T} + \varepsilon^2\mathsf I)
= (\mathsf G^{\mathsf T}\mathsf G + \varepsilon^2\mathsf I)\mathsf G^{\mathsf T}$;
multiplique à esquerda por $(\mathsf G^{\mathsf T}\mathsf G + \varepsilon^2
\mathsf I)^{-1}$ e à direita por $(\mathsf G\mathsf G^{\mathsf T} +
\varepsilon^2\mathsf I)^{-1}$, ambas invertíveis para $\varepsilon > 0$.
(b) Aqui $\mathsf G\mathsf G^{\mathsf T} = \|\mathbf g\|^2 = 2500$, um número,
logo $\mathsf R_\varepsilon = \mathsf G^{\mathsf T}\mathsf G/(2500 + \varepsilon^2)$,
de traço $2500/(2500 + \varepsilon^2)$. (c) O dado previsto é
$\mathbf g\cdot\mathbf m_\varepsilon = 0.50\times2500/(2500 + \varepsilon^2)$
e o resíduo $0.50\,\varepsilon^2/(2500 + \varepsilon^2)$. Para
$\varepsilon^2 = 0$: traço 1, resíduo 0; para 2500: traço 0.5, resíduo
\qty{0.25}{s}; para \num{10000}: traço 0.2, resíduo \qty{0.40}{s}. Mais amortecimento
fornece um modelo menor, uma resolução pior e um desajuste maior.
\end{solution}

\begin{solution}{exo:b3:seismology-imaging:12}
(a) Pelo exercício 5, $T = (2R/v)\sin(\Delta/2)$, logo $\mathrm dT/\mathrm d\Delta
= (R/v)\cos(\Delta/2)$. Uma corda que faz o ângulo $\Delta$ no centro é a
base de um triângulo isósceles com ângulo no vértice $\Delta$, então encontra o raio
na fonte em $i_0 = (\pi - \Delta)/2$, e $p = R\sin i_0/v =
(R/v)\cos(\Delta/2)$: o mesmo resultado. (b) O ponto de retorno é o ponto médio da
corda, à distância $R\cos(\Delta_1/2)$ do centro; equivalentemente
$r_1 = p_1v$. (c) Aqui $p(\Delta)/p_1 = \cos(\Delta/2)/\cos(\Delta_1/2)$, e
a fórmula dá $r_1 = R\exp\bigl(-\ln(1/\cos(\Delta_1/2))\bigr) =
R\cos(\Delta_1/2)$, como em (b); a velocidade é $r_1/p_1 = R\cos(\Delta_1/2)\big/
\bigl((R/v)\cos(\Delta_1/2)\bigr) = v$.
\end{solution}

\begin{solution}{pb:b3:seismology-imaging:1}
\textbf{1.} $i_0 = \arcsin 0.8 = \ang{53.13}$ e $p = 0.8/5.0 =
\qty{0.16}{s.km^{-1}}$.

\textbf{2.} O raio retorna onde $v = 1/p = \qty{6.25}{km.s^{-1}}$, ou seja, em
$z = (6.25 - 5.0)/0.05 = \qty{25}{km}$.

\textbf{3.} Derivando $\sin i = pv$ ao longo do raio, $\cos i\,\mathrm di/
\mathrm ds = pg\,\mathrm dz/\mathrm ds = pg\cos i$, logo $\mathrm di/\mathrm ds =
pg$: a direção gira a uma taxa constante, a curvatura de um círculo de raio
$1/(pg) = 1/(0.16\times0.05) = \qty{125}{km}$. No ponto de retorno o raio é
horizontal a \qty{25}{km} de profundidade e o centro situa-se verticalmente a \qty{125}{km}
acima dela: em $z = \qty{-100}{km}$, isto é, $v_0/g$ acima da superfície, onde a
velocidade extrapolada se anularia.

\textbf{4.} Na fonte, o raio do círculo é perpendicular ao raio sísmico,
logo forma o ângulo $i_0$ com a horizontal, e o centro fica a
$125\cos i_0$ de distância horizontal; por simetria,
$X = 2\times125\times\cos i_0 = 2\times125\times0.6 = \qty{150}{km}$, ou
$X = 2\sqrt{1 - p^2v_0^2}/(pg)$.

\textbf{5.} $\mathrm dT = \mathrm ds/v$ com $\mathrm ds = \mathrm di/(pg)$ e
$v = \sin i/p$: $\mathrm dT = \mathrm di/(g\sin i)$. Então
\[
  T = \frac2g\int_{i_0}^{\pi/2}\frac{\mathrm di}{\sin i} = \frac2g\ln\cot\frac{i_0}2 =
  \frac2g\operatorname{arccosh}\frac1{pv_0} .
\]
Aqui $\cot(\ang{26.57}) = 2$, logo
$T = 40\ln2 = \qty{27.73}{s}$.

\textbf{6.} Para $p = 0.15$: $pv_0 = 0.75$, $\sqrt{1 - 0.5625} = 0.6614$,
$X = 2\times0.6614/0.0075 = \qty{176.4}{km}$ e $T = 40\operatorname{arccosh}
(1.3333) = 40\times0.7954 = \qty{31.81}{s}$. Então $\Delta T/\Delta X =
4.08/26.4 = \qty{0.155}{s.km^{-1}}$, entre 0.15 e 0.16, como
$\mathrm dT/\mathrm dX = p$ exige.

\textbf{7.} $\mu_1 = 2700\times3500^2 = \qty{33.1}{GPa}$ e
$\mu_2 = 3300\times4500^2 = \qty{66.8}{GPa}$.

\textbf{8.} $k = 2\pi/(cT) = 2\pi/(3.99\times25) = \qty{0.06299}{rad.km^{-1}}$;
$q_1 = \sqrt{(3.99/3.5)^2 - 1} = 0.5474$; $q_2 = \sqrt{1 - (3.99/4.5)^2} = 0.4624$.

\textbf{9.} $\tan(kHq_1) = 66.825\times0.4624/(33.075\times0.5474) = 1.707$; o
modo fundamental toma $kHq_1 = \arctan1.707 = 1.041$, logo
$H = 1.041/(0.06299\times0.5474) = \qty{30.2}{km}$.

\textbf{10.} Em \qty{40}{s}: $k = 2\pi/170 = 0.03696$, $q_1 = 0.6888$,
$q_2 = 0.3287$, $\tan(kHq_1) = 0.964$, $kHq_1 = 0.7671$ e
$H = 0.7671/(0.03696\times0.6888) = \qty{30.1}{km}$. Os dois resultados concordam dentro da
precisão das velocidades de fase: a crosta tem cerca de \qty{30}{km} de espessura.

\textbf{11.} $\omega = 0.2513$ e $\qty{0.1571}{rad.s^{-1}}$, $k = 0.06299$ e
$\qty{0.03696}{rad.km^{-1}}$: $U \approx 0.0942/0.02603 = \qty{3.62}{km.s^{-1}}$,
menor que as duas velocidades de fase, como deve ser quando $c$ cresce com o
período.

\textbf{12.} $\lambda = 3.99\times25 = \qty{99.8}{km}$ e $4.25\times40 =
\qty{170}{km}$. O modo fundamental amostra até uma fração de seu
comprimento de onda: a onda de 25 s fica majoritariamente na crosta, a de 40 s
penetra bem no manto, então sua velocidade se aproxima de $\beta_2$. Dois períodos fixando o
mesmo $H$ é uma verificação de que o modelo de duas camadas é adequado.

\textbf{13.} Com as colunas das células 1 a 4 e as linhas A a F, em unidades
de \qty{1000}{km} ($\sqrt2\times100/1000 = 0.1414$):
\[
  \mathsf G = \begin{pmatrix}
  0.1 & 0.1 & 0 & 0 \\ 0 & 0 & 0.1 & 0.1 \\ 0.1 & 0 & 0.1 & 0 \\
  0 & 0.1 & 0 & 0.1 \\ 0.1414 & 0 & 0 & 0.1414 \\ 0 & 0.1414 & 0.1414 & 0
  \end{pmatrix}.
\]

\textbf{14.} $\mathsf C_d^{-1} = \operatorname{diag}(400, 400, 400, 400, 100,
100)$. Um raio de linha ou coluna contribui com $0.1^2\times400 = 4$, um raio
diagonal com $0.1414^2\times100 = 2$. A célula 1 é cruzada pelos raios A, C (4 cada) e E (2): 10. As
células 1 e 2 compartilham o raio A: 4; as células 1 e 3 compartilham C: 4; as células 1 e 4 compartilham E: 2; e
assim por diante por simetria, o que dá as quatro linhas enunciadas.

\textbf{15.} Um raio de linha ou coluna contribui com $0.1\times400\,\delta t = 40\,
\delta t$ para cada uma de suas células, um raio diagonal com $0.1414\times100\,\delta t =
14.14\,\delta t$. Célula 1: $40(0.16 - 0.36) + 14.14\times(-0.15) = -10.12$; célula
2: $40(0.16 + 0.11) + 14.14\times(-0.27) = 6.98$; célula 3: $40(-0.48 - 0.36) +
14.14\times(-0.27) = -37.42$; célula 4: $40(-0.48 + 0.11) + 14.14\times(-0.15) =
-16.92$.

\textbf{16.} Para $(1,1,1,1)$, cada linha de $\mathsf N$ soma 20. Para
$(1,-1,-1,1)$, a linha 1 dá $10 - 4 - 4 + 2 = 4$ e as demais seguem o mesmo
padrão. Para $(1,1,-1,-1)$, $10 + 4 - 4 - 2 = 8$, e para $(1,-1,1,-1)$,
$10 - 4 + 4 - 2 = 8$. Os vetores são ortogonais com norma 2, de modo que
$\mathsf N^{-1} = \sum_k \mathbf u_k\mathbf u_k^{\mathsf T}/(4\lambda_k)$; seu
elemento $(1,1)$ é $\tfrac14(\tfrac1{20} + \tfrac14 + \tfrac18 + \tfrac18) =
\tfrac{11}{80}$, seu elemento $(1,2)$ é $\tfrac14(\tfrac1{20} - \tfrac14 + \tfrac18
- \tfrac18) = -\tfrac4{80}$, seu elemento $(1,4)$ é $\tfrac14(\tfrac1{20} +
\tfrac14 - \tfrac18 - \tfrac18) = \tfrac1{80}$, e assim por diante.

\textbf{17.} $\hat{\mathbf m} = \mathsf N^{-1}\mathbf b$:
$\hat m_1 = (11\times(-10.121) - 4\times6.982 - 4\times(-37.418) - 16.921)/80
= -0.081$; analogamente $\hat m_2 = 1.844$, $\hat m_3 = -3.706$ e
$\hat m_4 = \qty{-0.931}{ms.km^{-1}}$.

\textbf{18.} $\delta v/v_0 = -v_0\,\delta s = -8.0\times10^{-3}\,\hat m$, em
por cento $-0.8\,\hat m$: célula 1 $+0.07$, célula 2 $-1.48$, célula 3 $+2.96$, célula 4
$\qty{+0.75}{\percent}$.

\textbf{19.} Todo elemento da diagonal de $\mathsf N^{-1}$ é $11/80$: cada célula tem
o erro padrão $\sqrt{0.1375} = \qty{0.371}{ms.km^{-1}}$, isto é,
$\qty{0.30}{\percent}$. A correlação das células 1 e 2 é $-4/11 = -0.36$:
elas compartilham o raio A, e uma lentidão atribuída a uma pode ser compensada pela outra.
As células 1 e 4, que compartilham apenas a diagonal E, menos precisa, são fracamente e
positivamente correlacionadas, $+1/11 = +0.09$.

\textbf{20.} Os resíduos previstos são $0.176$, $-0.464$, $-0.379$, $0.091$,
$-0.143$ e $\qty{-0.263}{s}$; os resíduos após o ajuste são $-0.016$,
$-0.016$, $0.019$, $0.019$, $-0.007$ e $\qty{-0.007}{s}$, e
$\chi^2 = 0.50$ para $6 - 4 = 2$ graus de liberdade: compatível com os erros
indicados; se houver algum desvio, ele tende a ser para o lado pequeno.

\textbf{21.} Os autovetores são os de $\mathsf N$ e os autovalores
$\lambda/(\lambda + 4)$: $20/24 = 5/6$, $4/8 = 1/2$, $8/12 = 2/3$ duas vezes. Como na
questão 16, $R_{11} = \tfrac14(\tfrac56 + \tfrac12 + \tfrac23 + \tfrac23) =
\tfrac23$; $R_{12} = \tfrac14(\tfrac56 - \tfrac12 + \tfrac23 - \tfrac23) =
\tfrac1{12}$; $R_{13} = \tfrac14(\tfrac56 - \tfrac12 - \tfrac23 + \tfrac23) =
\tfrac1{12}$; $R_{14} = \tfrac14(\tfrac56 + \tfrac12 - \tfrac23 - \tfrac23) = 0$.

\textbf{22.} Projeta-se $\mathbf b$ sobre os autovetores: $\mathbf b\cdot\mathbf u_k/4
= -14.37$, $0.849$, $12.80$ e $-9.40$ para os autovalores 20, 4, 8 e 8.
A estimativa amortecida divide-os por $\lambda + 4$ em vez de $\lambda$; para
a célula 3 (componentes $1, -1, -1, 1$),
$m_3 = -14.37/24 - 0.849/8 - 12.80/12 - 9.40/12 = \qty{-2.555}{ms.km^{-1}}$,
uma anomalia de velocidade de $\qty{+2.04}{\percent}$. O amortecimento puxa cada componente
para zero, sobretudo os mal restringidos.

\textbf{23.} O modelo não amortecido, que ajusta os dados dentro de seus erros
($\chi^2 = 0.50$ para 2 graus de liberdade), ao passo que o amortecido os desajusta
muito ($15.8$). O amortecimento só é necessário quando $\mathsf N$ é singular ou quase
isso; aqui seu menor autovalor, 4, não é pequeno, a estimativa é estável e
seus erros são moderados.

\textbf{24.} A linha 3 é $(1/12, 0, 2/3, 1/12)$: a estimativa amortecida da célula 3 é
dois terços de seu valor verdadeiro mais um doze avos das células 1 e 4, suas vizinhas
ao longo dos raios C e B, e nada da célula 2.

\textbf{25.} A célula 3, a inferior à esquerda, é rápida em cerca de
\qty{3}{\percent}; a célula 2 é lenta. Uma laje fria e densa de litosfera oceânica
afundando sob o continente é mais rápida que o manto circundante, de modo que a
célula rápida, profunda e do lado oceânico de K, é a laje; a célula lenta acima
dela, do outro lado, situa-se onde a cunha do manto é quente.

\textbf{26.} A partir das questões 18 e 19, a anomalia de velocidade P da célula da laje é
\textbf{\qty[text-series-to-math=true, uncertainty-mode=separate]{3.0 \pm 0.3}{\percent}}.
\end{solution}
